\documentclass[11pt]{article}
\usepackage[a4paper, margin=2.5cm]{geometry}
\usepackage{graphicx}
\usepackage{booktabs}
\usepackage{amsmath}
\usepackage[colorlinks=true, linkcolor=blue]{hyperref}
\usepackage{caption}
\captionsetup[figure]{name=Supplementary Figure, labelfont=bf}
\captionsetup[table]{name=Supplementary Table, labelfont=bf}

\graphicspath{{../figures/}}

\title{Supplementary Information for:\\
\large Conditional robustness of equatorial Kelvin waves in the real ocean from multi-mission altimetry and SWOT snapshots}
\author{Author One, Author Two}
\date{}

\begin{document}
\maketitle

%% ------------------------------------------------------------
\section*{Supplementary Note 1: Kelvin wave event detection and deduplication}

Candidate events were detected by ray tracing on the equatorial Hovm\"oller diagram (Methods), yielding 11 candidate rays. Ridge-intercept clustering on $\tau = t_\mathrm{start} - x_0/c$ merged co-ridge detections and produced 7 independent events (KE01--KE07), used as the sampling unit in all statistics.

\begin{figure}[h]
\centering
\includegraphics[width=0.95\textwidth]{p1_kelvin_events_detected.png}
\caption{Kelvin wave events detected from the DUACS equatorial Hovm\"oller diagram (2$^\circ$S--2$^\circ$N mean sea-level anomaly). Black lines mark detected rays (phase speed ${\sim}2.5$~m~s$^{-1}$) before ridge-intercept deduplication.}
\label{sfig:events}
\end{figure}

%% ------------------------------------------------------------
\section*{Supplementary Note 2: Westerly wind burst verification}

Wind forcing was verified with localized ERA5 10-m zonal wind analysis (Methods, ``Wind forcing verification''). Box-averaging over the full $150^\circ$E--$180^\circ$ search region yields easterly means for all seven events and would produce uniform false negatives; spatially localized detection is essential.

\begin{figure}[h]
\centering
\includegraphics[width=\textwidth]{p1_wwb_localized.png}
\caption{ERA5 equatorial ($2^\circ$S--$2^\circ$N) zonal wind around each Kelvin event onset. Red shading denotes westerlies; the dashed line marks the event onset; the star marks the westerly peak used in Supplementary Table~\ref{stab:wwb}.}
\label{sfig:wwb}
\end{figure}

\begin{table}[h]
\centering
\caption{Localized WWB detection per event (from \texttt{wwb\_event\_confirmation\_localized.json}). Peak westerly is the daily maximum of the $5^\circ$-smoothed, $2^\circ$S--$2^\circ$N-averaged zonal wind in the window from 20 days before to 5 days after onset. ``Edge'' indicates a peak at the $150^\circ$E search boundary; the signal may extend further west.}
\label{stab:wwb}
\begin{tabular}{lllrrrl}
\toprule
Event & Verdict & Peak day & Peak (m\,s$^{-1}$) & Lon ($^\circ$E) & Days $>3$~m\,s$^{-1}$ & Note \\
\midrule
KE01 & marginal  & 2023-01-03 & $+3.0$ & 150.0 & 0  & edge \\
KE02 & marginal  & 2023-02-22 & $+4.3$ & 150.0 & 2  & edge \\
KE03 & confirmed & 2023-05-19 & $+9.5$ & 150.0 & 6  & edge \\
KE04 & confirmed & 2023-07-19 & $+6.4$ & 157.2 & 9  & \\
KE05 & confirmed & 2023-07-31 & $+5.8$ & 154.5 & 8  & \\
KE06 & confirmed & 2023-10-07 & $+6.0$ & 163.2 & 13 & \\
KE07 & confirmed & 2023-12-17 & $+6.2$ & 166.8 & 11 & \\
\bottomrule
\end{tabular}
\end{table}

\clearpage

%% ------------------------------------------------------------
\section*{Supplementary Note 3: Scale dependence of the amplitude-based criterion}

The along-ray $\Lambda_\mathrm{min}$ uses the local maximum perturbation in a $\pm 1.5^\circ$ longitude $\times$ $\pm 3^\circ$ latitude window centred on the ray position, after $0.5^\circ$ boxcar smoothing. All three zones collapse to $\Lambda_\mathrm{min} \approx 1$, demonstrating that the amplitude criterion carries no zone information at the local-extreme scale (main text, Fig.~4b).

\begin{figure}[h]
\centering
\includegraphics[width=0.6\textwidth]{p4_lambda_along_ray.png}
\caption{Along-ray $\Lambda_\mathrm{min}$ versus zone-averaged $\Lambda$ for all event--zone pairs. The dashed diagonal marks equality; the dotted line marks $\Lambda_c \sim 1$.}
\label{sfig:alongray}
\end{figure}

%% ------------------------------------------------------------
\section*{Supplementary Note 4: $\Lambda_2$ correlation statistics and record-length caveat}

\begin{table}[h]
\centering
\caption{Correlations between $\log \Lambda_2$ and $\log$ amplitude ratio (from \texttt{lambda\_v2\_correlations.json}; quality filter $\mathrm{rms}_\mathrm{up} > 0.01$~m). The full sample is reported first; zone subsets follow. The Gilbert Islands retain only one valid pair after filtering and are omitted.}
\label{stab:corr}
\begin{tabular}{lrrrrr}
\toprule
Sample & $n$ & Spearman $\rho$ & $p$ & Pearson $r$ (log--log) & $p$ \\
\midrule
All event--zone pairs & 14 & $+0.420$ & 0.135 & $+0.394$ & 0.164 \\
TIW zone              & 7  & $+0.643$ & 0.119 & $+0.690$ & 0.086 \\
Line Islands          & 6  & $-0.543$ & 0.266 & $-0.634$ & 0.177 \\
\bottomrule
\end{tabular}
\end{table}

\begin{figure}[h]
\centering
\includegraphics[width=0.95\textwidth]{p4_lambda_v2.png}
\caption{$\Lambda_2$ by perturbation zone (left) and against the observed amplitude ratio (right), as in main-text Fig.~4c but showing all zones with the per-zone statistics of Supplementary Table~\ref{stab:corr}.}
\label{sfig:lambda2}
\end{figure}

\paragraph{Record-length caveat.} All event windows (50--65 days) are shorter than twice the longest resonant period ($2 \times 50 = 100$~days), so the 40--50~day frequency bin has limited resolution throughout the dataset. To assess sensitivity, we repeated the $\Lambda_2$ analysis with narrower resonant windows: 15--40~days and 15--30~days, which are fully resolved by all event records. The TIW-zone correlation sign and approximate magnitude are stable across window choices (see historical expansion results). Additionally, excluding the shortest-window event (KE05, 50 days), the TIW-zone Pearson $r = 0.63$ at $n = 6$ versus $r = 0.69$ at $n = 7$, and the rank correlation weakens modestly (Spearman $\rho = 0.43$ versus $0.64$).

\end{document}
